% Pista ana-23j2-q4 · Anàlisi
% Procedència: PAU Catalunya 2023 · Convocatòria de juny · Sèrie 5 · Qüestió 4
\documentclass[11pt,a4paper]{article}
\usepackage{pista}
\pistainfo{Catalunya 2023}{Convocatòria de juny}{Sèrie 5}

\begin{document}

\section*{\textcolor{blauPista}{Qüestió 4}}

\noindent En una carretera principal hi trobem el poble $A$. A 12 km del poble $A$ hi ha un encreuament $O$ amb una carretera secundària que talla perpendicularment la carretera principal. A 9 km de l'encreuament, a la carretera secundària, hi trobem el poble $B$. Es vol construir una torre de comunicacions $T$ en un punt de la carretera principal situat entre el poble $A$ i l'encreuament $O$. Aquesta torre ha d'estar connectada amb cadascun dels dos pobles en línia recta per cable. Sabem que instal·lar el cable entre la torre $T$ i el poble $B$ té un preu de $250\,\euro/\text{km}$ i, en canvi, instal·lar el cable entre la torre $T$ i el poble $A$ té un preu de $125\,\euro/\text{km}$. Determineu a quina distància de l'encreuament $O$ a la carretera principal cal situar la torre $T$ perquè el preu del cablejat sigui mínim i quin serà el valor d'aquest preu mínim. \hfill\textnormal{[2,5~punts]}

\pista{Comença per un dibuix i tria bé la variable: anomena $x$ la distància de l'encreuament $O$ a la torre $T$. Expressa en funció de $x$ la longitud de cada tram de cable (per al tram fins a $B$, fes servir Pitàgores) i munta la funció de cost $P(x)$ multiplicant cada longitud pel seu preu. A partir d'aquí és una optimització estàndard: deriva, iguala a zero i resol (en aïllar l'arrel hauràs d'elevar al quadrat). No oblidis comprovar amb el signe de $P'$ que el candidat és realment un mínim, i acaba avaluant $P$ en aquest punt per a donar el cost.}

\end{document}
